\section{怎样写比较复杂的说明文}

《中学语文教学大纲》对高中一年级学生提出“写比较复杂的”“说明的文章”的要求。所谓“比较复杂”，是指要说明的对象本身就比较复杂，因而说明的内容、结构和方法等随之也比较复杂起来。例如：《公共汽车》、《蜘蛛》、《故乡名胜》等，是要我们介绍一种比较简单的事物，而《古今交通工具》、《春蚕和蜘蛛》、《中国 --- 伟大的文明古国》等，则是要我们介绍比较复杂的事物了。即以《古今交通工具》而言，既要说明“古”的，又要说明“今”的；而且“交通工具”的内涵极广，包括海、陆、空甚至地下的交通工具都在说明的范围之内；何况交通工具的类别太多，天上的飞行物、水中的船只、陆上的车辆、地铁等等都在内；单以车辆而言，就有独轮手推车、自行车、三轮车、汽车、火车等等。至于“公共汽车”只不过是“汽车”中的一种罢了。很明显，随着说明内容的复杂化，我们的说明形式和方法也必须复杂起来。我们如果不掌握写作比较复杂的说明文的要领，文章必然会显得杂乱无章，顾此失彼，达不到“说明白了”的目的。

那么，我们写比较复杂的说明文时，应该注意哪些问题呢？

\textbf{1. 必须把精密观察与调查研究结合起来，以保证说明具有真实性和科学性。}

写比较简单的说明文，一般能仔细观察事物的外部特征就可以了；写比较复杂的说明文，则不能满足于把握事物的明显部分，还得特别留心观察事物的隐蔽部分，不但要把握事物静止时的特征，而且要特别留心观察事物变化中的特征。这些方面往往会因粗心大意而被忽略，而忽略了就不能真正了解事物。例如课文《蝉》，作者法布尔先观察蝉不等捕它的竿子触身就已飞去，这可算是视觉灵敏了吧？然而法布尔再观察竿子不发出一点声响伸向蝉时，蝉竟视而不见，即使竿子靠近了蝉的眼睛，它也不飞走，原来蝉竟然是个“瞎子”，它只凭身上发出的超声波感知事物，而不是凭眼睛看事物。再如观察蜘蛛结网，丝能粘住捕获的昆虫，而蜘蛛自己却可以自由来去，不会被粘住，这是为什么呢？再仔细观察一下，原来蜘蛛的脚底能分泌一种油脂性物质。如果不观察研究这些隐藏的、变化的特征，就不可能准确地说明这些事物。一篇说明文，不管你写得多么清楚明白，有条有理，只要观察不准确，内容不真实，不科学，也是不能称为较好的说明文的，因为说明文的内容必须要有准确的知识性，内容错了，形式再好，也是不行的。

有些事物的特征不可能全部观察得到。为了保证说明文的科学性，还要采用其他方式进行调查，搜集有关的可靠资料，充分注意到有关的专门知识，掌握那些准确反映事物特征的专门名词术语，再结合观察所得，进行分析研究。例如课文《向沙漠进军》，作者如果不调查研究我国西北干旱地区的地面径流和地下潜水等，不掌握充分的可靠的资料，是无法动笔的；再如课文《中国石拱桥》，作者如果不查阅有关记载中国各种石拱桥的资料，没有对石拱桥的特殊结构和原理作过切实的研究，也是无法动笔的。我们在学写比较复杂的说明文时，除了细致深入地观察客观事物外，还必须对被说明的对象深入调查研究，向内行人多请教，和同学们反复进行探讨，甚至查阅有关的书籍、材料，充分利用别人已经观察研究过的材料，这样才可能保证所写的文章具有科学性。

\textbf{2. 要根据被说明的事物自身的条理性安排好文章的结构。}

比较复杂的说明文不容易写好，主要在于被说明的事物有其复杂性。既然这样，我们在进行说明时，首先就给读者总说一下这个事物的概况，让人有一个总的印象；然后再具体介绍事物的各个组成部分或历史发展的情况，最后再回到总体上来，说明对事物本质的认识，指出写作本文的目的。这样的结构方式，既能抓住要点，纲举目张，又能有条有理，说得明白。例如课文《人民英雄永垂不朽》，首先点明人民英雄纪念碑的位置和建碑的起因。“它象征着先烈们的丰功伟绩，标志着全国人民对先烈的怀念。”接着按照瞻仰纪念碑的顺序依次说明：先概括地从白玉台阶、二层平台写到碑身四周的汉白玉栏杆，再突出介绍碑身正面毛泽东同志的题字和碑身背面周恩来同志书写的碑文；又从东面、南面写到西面，按照历史顺序，详细说明碑身上的一幅幅浮雕；最后又回到碑身的正面（北面），集中写十幅浮雕中最大的一幅 --- “胜利渡长江，解放全中国”。文章以瞻仰纪念碑后的感想结尾，点明中心，揭示本质特征。这样写，主次分明，条理井然，能给人一个清楚、完整的印象。

我们再谈谈怎样具体介绍事物的各个组成部分，因为这是说明文的主干部分，处理得好就眉目清楚，处理得不好就杂乱无章。要想处理得好，必须善于发现事物、事理自身的条理性。这种条理性是客观存在的。例如人民英雄纪念碑有碑座、碑身以及碑身、四面的浮雕这样一个比较合理的顺序；再如水稻的植株自下而上由根、茎、叶、穗四部分组成；保温瓶自外而内由瓶壳和瓶胆组成；印一本书，必须先拣字排版，再上印刷机，然后才能折页装订；种庄稼，一定要先选种，再播种，接着要施肥除草，最后才是收获；在现代自然科学中，基础学科是应用学科的基础，而基础学科中的物理和数学又是化学、天文学、地学、生物学这四门学科的基础······可以说，一切事物、事理都有自己的规律性、条理性。说明文的结构，至少它的主干部分，正是按照这种规律性、条理性来安排的。李四光不就是依据人类发展进程所呈现出来的条理性，把《人类的出现》的主干部分分作“古猿 --- 猿人 --- 古人 --- 新人”四层的么？假如我们也能象他那样掌握住说明对象固有的条理性，那么，安排文章结构就不是多难的事了。

应该注意：复杂事物本身的几个组成部分或几个方面是有机联系在一起的，文章中相应的几个段落也应象接榫那样紧密相联；同时，每个段落的首句还应交代所要说明的具体范围或内容，总领全段。这样，就容易使全文中有几个横向说明的部分或纵向说明的部分既衔接自然，又头绪清晰。

\textbf{3. 要善于综合运用多种说明方法。}

要把比较复杂的事物说明清楚，让读者具体、明白地了解说明对象的表面特征和本质特征，就要掌握多种说明方法，并善于综合运用。说明的方法主要有：

（1）提问题和下定义。

写说明文，常常把自己所要说明的中心问题或几个方面的问题提出来，以便引起读者的注意，启发读者思考，同时使文章纲目清楚。提问题的位置，一般放在一篇文章的开头，有时也可以用在文章中间的关键地方，起到强调、提示和引导的作用。问题提出来了，就得针对问题进行回答，不能问而不答，或者答非所问。回答的方式，往往是用简要的话，把事物的性质和所包含的意义概括出来，使读者形成一个比较明确的概念，这就是通常所说的下定义。定义的语言必须准确、简炼，概括的内容必须正确、肯定。

例如：提问题“什么是‘机器人’？”课文里给它的定义是“在电子计算机的控制下，能够模仿人的某些动作和智能，成为人们从事生产劳动和科学实验的助手”。这个定义有确定的意义：它不同于人，因为它是在电子计算机的控制下工作的；它又不同于一般机器，因为它能够模仿人的某种动作和智能。从这个意义上看，它只能是“机器人”，而不是别的什么。由于各种事物都有其主要的特征和本质，因此，对事物下定义，应该抓住其主要的方面。如有的定义着重说明事物的特征，有的着重说明事物的作用，有的定义则从几个方面来说明，既揭示事物的特征和本质，又说明了作用。当然以最后一种下定义的方法表述最严密，最有科学性。这里有个原则，就是不论从哪个角度下定义，都应该能使读者了解这种事物同别的事物的区别，了解你要说明的这种事物的基本特征。

（2）分类和举例。
定义只能对事物的本质作概括的说明，而要了解事物的具体内容，还得作进一步的分析说明。这种分析说明往往采取分类的方法。给事物分类，首先要熟悉事物的特征；其次要按一定的标准，将有共同特点的事物归纳在一起，然后把每类事物说得具体、清楚。这种分类法，我们已在《怎样写比较简单的说明文》中作了介绍，不再赘述。

问题是，有些事物或事理，不容易说得清楚明白，这时往往需要举例说明。如在课文《奇特的激光》中，就有许多地方运用了举例的方法。其中有一段写了“激光的方向性也有很多的功能。”什么是“方向性”呢？又有哪些功能呢？文章中举例说明道：“如我们挖掘长距离的隧道或矿山巷道，可以用激光束来导向，这样挖出来的地道又准又直。建设高层楼房，可以用激光束来代替吊线，这样建成的楼房保证上下笔直。”这个例子有力地说明了激光在方向性方面的功能，使读者对于激光的“奇特”，有了清楚的认识，增长了知识。需要注意：举例要有代表性，真实可靠，能恰当地说明事物的特性；同时要力求浅显、具体，有趣味性，以便增强说明的效果。

（3）比较和引用。

事物的特征往往在比较中显现出来，有比较才有鉴别。如课文《宇宙里有些什么》这样写道：“许多红色的星星很大很大，最大的装得下八十万万个太阳。这些星星是由非常稀薄的气体状态的物质组成的。最稀薄的，密度只有地球上空气的几万分之一，比我们用抽气机造成的真空还要稀薄得多。”这里，为了说明恒星之大，把它同太阳比较；为了说明有些恒星密度的稀薄，把它同抽气机造成的“真空”作比较。这样就能更清楚地说明恒星的特性。应该注意：拿来相比的事物，只能在某一相同的角度或范畴进行比较，从而看出它们不同的特征；而且要尽可能选择为大家所熟悉的、容易理解的事物来进行比较，使读者一看就懂。如果不是在某一相同角度或范畴，事物之间就无从比较；如果用来作比的事物比原来的事物还陌生、还难懂，那就不能达到说明事物的目的了。

为了有力地说明事物，有时还可以引用一些别人的言论或有关的资料，使读者更加明白事物的特征和本质。如《中国石拱桥》一文，为了说明赵州桥施工技术的“巧妙绝伦”，引用唐朝宰相张嘉贞的话：“创造奇特，人不知其所以为。”为了说明赵州桥闻名于世，引用意大利著名旅行家马可·波罗的话：“是世界上独一无二的。”为了说明赵州桥和芦沟桥的来龙去脉，文中还大量引用了有关的史料和建筑资料。应该注意：引用名言和资料，必须准确可靠；还必须从实际需要出发，只有引来能够确切说明问题的才引用，不能毫无选择地堆砌。

（4）举数字和列图表。

有些事物可以从数量上表明它的特征和本质，往往需要运用一些数字。举数字有时更便于说明问题。如课文《向沙漠进军》为了说明沙漠地区日照时间特别长，举了这样的数字：沙漠地区“一年达到三千小时，而长江流域只有一千五百小时，华北地区也不过二千五百小时。”这样写，就很具体明确，使读者一看就明白。又如课文《石油的用途》一文中说：“燃烧一公斤石油可以产生一万大卡的热量，而燃烧一公斤煤只有五千到六千五百大卡的热量。”两组数字的对比，说明了石油可燃性好、发热量高这个特征。应该注意：说明事物特征的数字必须准确无误，确定的数字必须一丝不差，估计的数字则应力求近似。为了达到说明的准确，写进文章里的数字必须经过核实才行。

对于有些比较复杂的事物，我们可以列出图表来帮助说明。例如课文《石油的用途》的课后“思考和练习”第三题，附了一份《豆腐粉食用说明书》，其中有一份“配料表”。我们如果对各种豆制品的配料不了解，看了这份“配料表”，对于嫩豆腐、老豆腐、豆浆各配什么料，该配多少料，就一清二楚了。除了表格，有时也用插图或照片来帮助说明某一事物，使说明更加清楚、醒目，而且直观、形象，使读者易于认识。

总之，在写比较复杂的说明文时，应根据说明的需要，灵活地综合运用各种说明方法，有助于把事物的表面特征和本质特征说明得一清二楚，具体恰当。

4.适当地运用描写和议论，增强文章的形象性、深刻性和可读性。说明文是以“说明”为主要表达方式的，但是并不完全排斥描写和议论。即使以说明《这下面没有水，再换个地方挖》那幅漫画内容的简单说明文而言，如果我们能扣住画题，用平实、质朴的语言，条理分明地介绍画面上的水、土、井、人和画题五部分内容，而在说明“人”的表面特征时，用一两句话描画一下挖井人的神态，并在结尾处，用一句议论的话点明作者是在对那种做事不专一、无恒心的人进行讽刺，这样不就既符合写“说明性文字”的要求，又恰到好处地综合运用了描写、叙述、议论等表达方式了吗？至于比较复杂的说明文，在说明为主的前提下，往往要穿插一些描写和议论。例如沈括《雁荡山》一文中的“予观雁荡诸峰，皆峭拔险怪，上耸千尺，穷崖巨谷，不类他山。自岭外望之，都无所见；至谷中则森然干霄”一节，就是颇为精彩的描写；再如李四光《人类的出现》的最末一节，则纯是议论了。这说明只要运用恰当，文字简洁就行。

在说明文中运用一些描写手法，不仅能使文章生动具体，并能借助形象把深奥的道理说得浅显易懂，易为读者所理解。特别在文艺性说明文中，它还往往通过完整的故事、具体生动的描绘，拟人比喻等修辞手法，把事物或事理说得深入浅出，使文章活泼有趣，不仅给人以知识，而且还给人以艺术的享受。例如布封的《琥珀》、《松鼠》，法布尔的《蝉》，伊林的《天气隆下》等文，无不是脍炙人口的文艺性说明文。这些作者以其敏锐、深入的观察，准确、生动的文艺笔调，向读者兴味盎然地介绍各种科学知识，赢得了全世界儿童的喜爱，这正说明了他们的表达手法的成功。如果我们要写文艺性说明文，就该向这些大师们好好学习。

在说明文中运用一些议论手法，则能让读者对所说明的道理了解得更加深刻，有时则能让读者更清楚地认识到文章的现实意义。例如课文《雄伟的天安门》，在开头叙述北京的建都，筑城等情况后，紧接着就是一段议论：“封建帝王最害怕人民，他们妄想用重重城墙来保护他们自己。但一旦革命爆发，这些城墙也只是纸壁墙，并不能使他们免于灭亡。”这几句议论，揭示了城墙的本质，使说明更加透彻；如果删去了这几句议论，读者就只能从表面现象来认识城墙的意义了。

比较复杂的说明文，往往是说明、描写、议论这三种主要表达方式交叉使用，互为补充的。但是，应该明确：运用描写和议论，一定要服从于说明的需要。不论是一段描写还是几句议论，都要成为全篇说明文的有机组成部分，决不能离开说明的中心而无目的地描写或议论；同时，除了特殊类型的文艺性说明文外，一般说明文不能有过多的描写或议论，否则很容易失去说明文的特点，把说明文写成了记叙文或者议论文。
\newpage